\documentclass[11pt,a4paper]{article}

% ── Packages ─────────────────────────────────────────────────────────────────
\usepackage[utf8]{inputenc}
\usepackage[T1]{fontenc}
\usepackage[french]{babel}
\usepackage{geometry}
\geometry{top=2.5cm,bottom=2.5cm,left=2.5cm,right=2.5cm}
\usepackage{lmodern}
\usepackage{microtype}
\usepackage{graphicx}
\usepackage{booktabs}
\usepackage{tabularx}
\usepackage{array}
\usepackage{xcolor}
\usepackage{colortbl}
\usepackage{fancyhdr}
\usepackage{titlesec}
\usepackage{caption}
\usepackage{float}
\usepackage{amsmath}
\usepackage{amssymb}
\usepackage{hyperref}
\usepackage{enumitem}
\usepackage{tcolorbox}
\usepackage{parskip}

% ── Couleurs ──────────────────────────────────────────────────────────────────
\definecolor{meteoBlue}{HTML}{1D3557}
\definecolor{meteoRed}{HTML}{E63946}
\definecolor{meteoGray}{HTML}{F1FAEE}
\definecolor{meteoCyan}{HTML}{457B9D}

% ── Mise en page ─────────────────────────────────────────────────────────────
\pagestyle{fancy}
\fancyhf{}
\fancyhead[L]{\small\color{meteoBlue}\textbf{Météorage} — Phase 2bis : Amélioration du Neural Hawkes}
\fancyhead[R]{\small\color{meteoBlue}Rapport technique}
\fancyfoot[C]{\small\color{gray}\thepage}
\renewcommand{\headrulewidth}{0.4pt}

% ── Titres ────────────────────────────────────────────────────────────────────
\titleformat{\section}[block]
  {\Large\bfseries\color{meteoBlue}}
  {\thesection}{1em}{}[\vspace{-0.3em}\rule{\linewidth}{0.4pt}]
\titleformat{\subsection}[block]
  {\large\bfseries\color{meteoCyan}}
  {\thesubsection}{1em}{}
\titleformat{\subsubsection}[block]
  {\normalsize\bfseries\color{meteoBlue}}
  {\thesubsubsection}{1em}{}

\captionsetup{font=small, labelfont=bf, margin=10pt}

\newenvironment{keybox}[1]
  {\begin{tcolorbox}[colback=meteoGray,colframe=meteoBlue,title={\textbf{#1}},
    fonttitle=\small, boxrule=0.6pt, arc=3pt]}
  {\end{tcolorbox}}

\hypersetup{colorlinks=true, linkcolor=meteoBlue, urlcolor=meteoCyan, citecolor=meteoRed}

% ─────────────────────────────────────────────────────────────────────────────
\begin{document}

% ── Page de titre ─────────────────────────────────────────────────────────────
\begin{titlepage}
  \pagecolor{meteoBlue}
  \color{white}
  \vspace*{2cm}
  \begin{center}
    {\Huge\bfseries Prédiction de fin d'orage\\[0.4em]}
    {\large\bfseries Phase 2bis — Amélioration du Neural Hawkes Process}\\[2em]
    \rule{0.5\linewidth}{0.8pt}\\[1.5em]
    {\large Réponse à l'appel d'offre \textbf{Météorage}}\\[0.5em]
    {\normalsize Enrichissement des features, architecture Transformer,\\
    sortie gaussienne en espace log, et \textbf{Mixture of Experts Spatial}}\\[3em]
    {\normalsize Mars 2026}
  \end{center}
  \vfill
  \begin{center}
    \small\textit{4 itérations $\cdot$ 7 variantes $\cdot$ 1\,682 sessions d'alerte $\cdot$ 20 features}
  \end{center}
\end{titlepage}
\nopagecolor

\tableofcontents
\newpage

% ═══════════════════════════════════════════════════════════════════════════════
\section{Rappel : le processus de Hawkes pour la modélisation d'orages}
% ═══════════════════════════════════════════════════════════════════════════════

\begin{keybox}{Intuition}
Un \textbf{processus de Hawkes} modélise un phénomène courant dans les orages : \emph{un éclair augmente la probabilité qu'un autre éclair survienne peu après}. L'activité croît en cascade puis retombe progressivement.
\end{keybox}

Le taux d'arrivée (intensité) des éclairs est modélisé comme :
\begin{equation}
\lambda(t) = \underbrace{\mu}_{\text{taux de fond}} + \underbrace{\sum_{t_i < t} \alpha \cdot \beta \cdot e^{-\beta (t - t_i)}}_{\text{excitation par les éclairs passés}}
\end{equation}

\begin{itemize}[nosep]
  \item $\mu$ : taux de fond — probabilité ``naturelle'' d'un éclair, même sans activité récente
  \item $\alpha$ : force d'excitation — combien chaque éclair augmente le taux d'arrivée
  \item $\beta$ : taux de décroissance — vitesse à laquelle l'effet d'un éclair s'estompe
\end{itemize}

\textbf{Prédiction de fin d'orage} : quand $\lambda(t)$ redescend vers $\mu$, l'orage est considéré terminé.

\subsection{Du Hawkes classique au Neural Hawkes}

Le Hawkes classique suppose une décroissance exponentielle fixe, trop restrictive pour les dynamiques complexes d'un orage. Le \textbf{Neural Hawkes} remplace la fonction de déclenchement par un réseau de neurones qui apprend la dynamique temporelle directement.

\subsection{Modèle de Phase 2 (référence)}

\begin{keybox}{Performance Phase 2 — GRU v1}
\begin{tabular}{@{}ll@{}}
Architecture : & GRU, $h=32$, 4 features (icloud, amplitude, distance, maxis) \\
MAE = 22.11 min & RMSE = 32.50 min \quad Biais = $-$1.61 min \quad Médiane = 18.00 min
\end{tabular}
\end{keybox}

\noindent Pour référence, les autres modèles de Phase 2 étaient nettement moins performants : XGBoost AFT (MAE~=~57.28), Baseline 30~min (MAE~=~60.26), BNN MC~Dropout (MAE~=~68.20), Hawkes classique (MAE~=~83.99).


% ═══════════════════════════════════════════════════════════════════════════════
\section{Axe 1 : Enrichissement des features}
% ═══════════════════════════════════════════════════════════════════════════════

\subsection{Motivation}

Le modèle de Phase 2 n'utilise que 4 features brutes par événement. Plusieurs indicateurs de dissipation ne sont pas exploités :

\begin{itemize}[nosep]
  \item \textbf{Direction} : un orage qui s'éloigne de l'aéroport a plus de chances de se terminer.
  \item \textbf{Densité temporelle} : une baisse du nombre d'éclairs/minute signale la dissipation.
  \item \textbf{Ratio IC/CG} : les éclairs intra-nuage (IC) précèdent souvent les CG ; un ratio IC croissant peut signaler la fin.
\end{itemize}

\subsection{Features ajoutées (8 nouvelles, 12 au total)}

\begin{table}[H]
\centering
\small
\begin{tabular}{@{}lll@{}}
\toprule
\textbf{\#} & \textbf{Feature} & \textbf{Description} \\
\midrule
0--3  & icloud, amplitude, dist, maxis & Features de Phase 2 (inchangées) \\
4--5  & azimuth\_sin, azimuth\_cos     & Direction encodée cycliquement \\
6     & dt\_since\_last                & Temps inter-événement (normalisé) \\
7     & dt\_since\_last\_cg            & Temps depuis le dernier CG \\
8     & density\_5min                  & Nombre d'éclairs dans les 5 dernières minutes \\
9     & ic\_ratio\_recent              & Ratio IC/total dans les 10 derniers éclairs \\
10    & dist\_to\_centroid             & Distance au centroïde mobile des 10 derniers \\
11    & dist\_trend                    & Tendance de la distance (positif = éloignement) \\
\bottomrule
\end{tabular}
\caption{Les 12 features par événement dans le modèle enrichi.}
\end{table}


% ═══════════════════════════════════════════════════════════════════════════════
\section{Axe 2 : Architecture Transformer avec attention temporelle}
% ═══════════════════════════════════════════════════════════════════════════════

\subsection{Pourquoi remplacer le GRU ?}

Le GRU traite la séquence de manière \emph{séquentielle} : l'information doit traverser chaque pas de temps. Pour les longues sessions (100+ éclairs), les dépendances à longue portée sont diluées.

Le \textbf{Transformer} utilise un mécanisme d'\textbf{attention} où chaque position accède directement à toutes les positions passées :
\begin{equation}
\text{Attention}(Q, K, V) = \text{softmax}\left(\frac{QK^\top}{\sqrt{d_k}}\right) V
\end{equation}
avec un \textbf{masque causal} empêchant chaque position de voir le futur.

\subsection{Encodage positionnel continu}

Contrairement au Transformer classique (indices discrets), notre modèle utilise les \textbf{timestamps réels} (en minutes) :
\begin{equation}
\text{PE}(t, 2i) = \sin\left(\frac{t}{10000^{2i/d}}\right), \quad \text{PE}(t, 2i+1) = \cos\left(\frac{t}{10000^{2i/d}}\right)
\end{equation}

\begin{keybox}{Pourquoi l'encodage continu ?}
L'espacement temporel entre éclairs est physiquement informatif. Un intervalle de 10~min vs 30~s a une signification très différente. L'encodage continu préserve cette information, là où un encodage discret ne verrait que ``position $n$ vs $n+1$''.
\end{keybox}

\subsection{Architectures testées}

\begin{table}[H]
\centering
\small
\begin{tabular}{@{}lccccc@{}}
\toprule
\textbf{Variante} & $d_{model}$ & \textbf{Têtes} & \textbf{Couches} & \textbf{FFN} & \textbf{Paramètres} \\
\midrule
Transformer       & 64  & 4 & 3 & 128 & $\sim$80K \\
Transformer Large & 128 & 8 & 4 & 256 & $\sim$400K \\
Par aéroport      & 64  & 4 & 2 & 128 & 5 $\times$ 60K \\
\bottomrule
\end{tabular}
\caption{Configurations des architectures Transformer testées.}
\end{table}


% ═══════════════════════════════════════════════════════════════════════════════
\section{Axe 3 : Architecture multi-tâche et sortie gaussienne}
% ═══════════════════════════════════════════════════════════════════════════════

\subsection{Le problème de l'extrapolation d'intensité}

Le modèle de Phase 2 prédisait le temps restant en \emph{extrapolant} manuellement la décroissance de $\lambda(t)$. Pendant une tempête active, les intensités récentes \emph{augmentent}, donnant un taux de décroissance quasi nul. Le modèle prédit alors que l'orage va durer indéfiniment.

\begin{keybox}{Constat — Itération 0 (extrapolation brute)}
Toutes les variantes avec features enrichies et/ou Transformer produisent une MAE~$\approx$~50 min et un biais de +30 min. Le problème ne vient pas du modèle, mais du mécanisme de prédiction.
\end{keybox}

\subsection{Itération 1 : tête de régression directe}

Au lieu d'extrapoler l'intensité, on ajoute une \textbf{seconde tête de sortie} qui prédit directement le temps restant depuis le hidden state :
\begin{equation}
\hat{y}_t = \text{Softplus}\big(\text{MLP}(h_t)\big) \geq 0
\end{equation}

La perte combinée est :
\begin{equation}
\mathcal{L} = \underbrace{\mathcal{L}_{\text{NLL}}}_{\text{processus ponctuel}} + \lambda \cdot \underbrace{\mathcal{L}_{\text{Huber}}(\hat{y}, y)}_{\text{régression directe}}
\end{equation}

\begin{keybox}{Résultat — Itération 1}
Biais réduit de +30 à +11 min, MAE de 50 à 27 min. Progrès significatif mais biais résiduel encore trop élevé.
\end{keybox}

\subsection{Itération 2 : Gaussian NLL en espace logarithmique}

Le temps restant $y$ a une distribution très \textbf{asymétrique} (beaucoup de valeurs proches de 0 en fin de session). La transformation $z = \log(1 + y)$ normalise cette distribution.

Le modèle prédit $(\mu, \log\sigma)$ en espace log, avec une \textbf{NLL gaussienne} :
\begin{equation}
\mathcal{L}_{\text{Gauss}} = \frac{1}{2} \left[ \frac{(z - \mu)^2}{\sigma^2} + 2\log\sigma + \log(2\pi) \right]
\end{equation}

La prédiction finale en minutes est $\hat{y} = e^{\mu} - 1$, et l'incertitude est naturellement calibrée.

\begin{keybox}{Résultat — Itération 2 (finale)}
\textbf{MAE = 21.33 min} (vs 22.11 pour Phase 2, $-$3.5\%), \textbf{biais = +1.96 min}, calibration 2$\sigma$ = 96.4\%.
\end{keybox}


% ═══════════════════════════════════════════════════════════════════════════════
\section{Axe 4 : Estimation d'incertitude bayésienne}
% ═══════════════════════════════════════════════════════════════════════════════

\subsection{MC Dropout}

Le \textbf{Monte Carlo Dropout} maintient le dropout actif pendant l'inférence et effectue $N=50$ passages forward. La variance des prédictions capture l'\textbf{incertitude épistémique}.

\subsection{Couches variationnelles}

Les poids sont des \textbf{distributions gaussiennes} $q(\mathbf{w}) = \mathcal{N}(\mu_w, \sigma_w^2)$ au lieu de valeurs fixes. L'optimisation minimise l'ELBO :
\begin{equation}
\mathcal{L}_{\text{ELBO}} = \mathcal{L}_{\text{NLL}} + \mathcal{L}_{\text{Reg}} + \beta \cdot \text{KL}\big(q(\mathbf{w}) \| p(\mathbf{w})\big)
\end{equation}

\subsection{Résultats bayésiens}

\begin{table}[H]
\centering
\small
\begin{tabular}{@{}lcccc@{}}
\toprule
\textbf{Méthode} & \textbf{MAE} & \textbf{Biais} & \textbf{1$\sigma$ cal.} & \textbf{2$\sigma$ cal.} \\
\midrule
MC Dropout (multi-tâche)  & 28.79 & +15.54 & 10.9\% & 26.9\% \\
Variational (multi-tâche) & 28.56 & +14.13 & 11.7\% & 27.0\% \\
\midrule
\textbf{GRU Gaussian (log-space)} & \textbf{21.33} & \textbf{+1.96} & \textbf{90.9\%} & \textbf{96.4\%} \\
\bottomrule
\end{tabular}
\caption{Comparaison des approches bayésiennes. Le passage en Gaussian NLL log-space résout à la fois le biais et la calibration.}
\end{table}

L'incertitude du GRU Gaussian est \textbf{sur-calibrée} (91\% dans 1$\sigma$ vs idéal 68\%). Le modèle est ``trop prudent'' — préférable dans un contexte opérationnel.


% ═══════════════════════════════════════════════════════════════════════════════
\section{Axe 5 : Modèles spécialisés par aéroport}
% ═══════════════════════════════════════════════════════════════════════════════

\begin{table}[H]
\centering
\small
\begin{tabular}{@{}lccrcc@{}}
\toprule
\textbf{Aéroport} & \textbf{Train} & \textbf{Test} & \textbf{MAE locale} & \textbf{MAE globale} & $\Delta$ \\
\midrule
Ajaccio  & 268 & 59  & \textbf{19.62} & 21.47 & $-$1.85 \\
Bastia   & 292 & 65  & \textbf{21.14} & 21.77 & $-$0.63 \\
Biarritz & 306 & 75  & 23.42 & \textbf{22.74} & +0.68 \\
Nantes   & 106 & 27  & \textbf{20.92} & 21.00 & $-$0.08 \\
Pise     & 373 & 111 & 22.85 & \textbf{20.17} & +2.68 \\
\midrule
\textbf{Ensemble} & & & 21.93 & \textbf{21.33} & +0.60 \\
\bottomrule
\end{tabular}
\caption{Modèle local vs global (v3 Gaussian, log-space). L'ensemble obtient un biais de +0.76 min.}
\end{table}

Les modèles locaux améliorent Ajaccio ($-$1.85 min) mais dégradent Pise (+2.68 min, moins de données per-class pour le modèle local). Le modèle global est légèrement meilleur en agrégé.


% ═══════════════════════════════════════════════════════════════════════════════
\section{Résultats comparatifs}
% ═══════════════════════════════════════════════════════════════════════════════

\begin{table}[H]
\centering
\small
\renewcommand{\arraystretch}{1.3}
\begin{tabular}{@{}llrrrrr@{}}
\toprule
\textbf{Iter.} & \textbf{Variante} & \textbf{MAE} & \textbf{RMSE} & \textbf{Méd.} & \textbf{Biais} & \textbf{P90} \\
\midrule
Réf. & Baseline 30 min       & 60.26 & ---   & ---   & $-$52.32 & --- \\
Réf. & XGBoost AFT           & 57.28 & 90.69 & 33.18 & $-$26.74    & 149.82 \\
Réf. & GRU v1 (Phase 2)      & 22.11 & 32.50 & 18.00 & $-$1.61  & 40.12 \\
\midrule
v2   & GRU multi-tâche       & 26.59 & 33.67 & 24.18 & +11.04   & 45.05 \\
v2   & Transformer multi-tâche & 29.17 & 35.55 & 28.64 & +14.10 & 44.86 \\
\midrule
\rowcolor{meteoGray}
\textbf{v3} & \textbf{GRU Gaussian (log)} & \textbf{21.33} & \textbf{30.25} & \textbf{15.96} & \textbf{+1.96} & \textbf{42.13} \\
v3   & Transformer Gaussian  & 23.25 & 31.59 & 19.34 & +4.23    & 43.12 \\
v3   & TF Large Gaussian     & 24.31 & 33.90 & 19.19 & +4.22    & 47.01 \\
v3   & Ensemble par aéroport & 21.93 & 31.88 & 16.22 & +0.76    & 43.37 \\
\bottomrule
\end{tabular}
\caption{Comparaison de toutes les variantes testées. Erreurs en minutes. La ligne surlignée est le meilleur modèle.}
\end{table}

\begin{figure}[H]
  \centering
  \includegraphics[width=0.95\linewidth]{figures/phase2bis_ablation.pdf}
  \caption{Comparaison ablative : MAE et biais de toutes les variantes testées.}
\end{figure}


% ═══════════════════════════════════════════════════════════════════════════════
\section{Analyse détaillée du meilleur modèle}
% ═══════════════════════════════════════════════════════════════════════════════

\begin{keybox}{Meilleur modèle — GRU Gaussian (log)}
MAE = 21.33 min \quad RMSE = 30.25 min \quad Médiane = 15.96 min\\
Biais = +1.96 min \quad Calibration 2$\sigma$ = 96.4\%\\
Homogénéité entre aéroports : écart max de 2.57 min sur la MAE.
\end{keybox}

\subsection{Performance par aéroport}

\begin{figure}[H]
  \centering
  \includegraphics[width=0.85\linewidth]{figures/phase2bis_per_airport.pdf}
  \caption{Performance par aéroport — modèle global vs local.}
\end{figure}

\subsection{Prédiction vs réalité}

\begin{figure}[H]
  \centering
  \includegraphics[width=0.65\linewidth]{figures/phase2bis_scatter.pdf}
  \caption{Scatter plot du GRU Gaussian. La diagonale représente la prédiction parfaite.}
\end{figure}

\subsection{Distribution des erreurs}

\begin{figure}[H]
  \centering
  \includegraphics[width=0.85\linewidth]{figures/phase2bis_error_dist.pdf}
  \caption{Distribution des erreurs (pred $-$ true) pour les trois variantes Gaussian.}
\end{figure}

\subsection{Performance selon l'avancement de la session}

\begin{figure}[H]
  \centering
  \includegraphics[width=0.85\linewidth]{figures/phase2bis_mae_by_fraction.pdf}
  \caption{MAE et biais en fonction de la progression dans la session d'orage.}
\end{figure}

\subsection{Calibration de l'incertitude}

\begin{figure}[H]
  \centering
  \includegraphics[width=0.6\linewidth]{figures/phase2bis_calibration.pdf}
  \caption{Diagramme de fiabilité. Courbe au-dessus de la diagonale = incertitude conservatrice (sur-couverture).}
\end{figure}


% ═══════════════════════════════════════════════════════════════════════════════
\section{Discussion et enseignements}
% ═══════════════════════════════════════════════════════════════════════════════

\subsection{Pourquoi le GRU bat le Transformer ?}

Contrairement aux attentes, le GRU surpasse systématiquement le Transformer. Trois hypothèses :

\begin{enumerate}[nosep]
  \item \textbf{Taille du dataset} : 1\,345 sessions d'entraînement, insuffisant pour un Transformer qui a besoin de beaucoup de données pour apprendre des patterns d'attention efficaces.
  \item \textbf{Longueur des séquences} : les sessions font en moyenne 30--50 éclairs. Le GRU est très efficace sur des séquences courtes.
  \item \textbf{Inductive bias} : la récurrence du GRU encode naturellement la ``décroissance temporelle'', ce qui correspond exactement au processus de Hawkes.
\end{enumerate}

\subsection{Impact du log-target}

Le passage Huber loss $\rightarrow$ Gaussian NLL en espace $\log(1+t)$ a été la plus grande amélioration :
\begin{itemize}[nosep]
  \item Biais : $+11.04 \to +1.96$ min ($-82\%$)
  \item MAE : $26.59 \to 21.33$ min ($-20\%$)
  \item Médiane : $24.18 \to 15.96$ min ($-34\%$)
\end{itemize}

La distribution asymétrique du temps restant rend la régression en espace linéaire biaisée vers les grandes valeurs.


% ═══════════════════════════════════════════════════════════════════════════════
\section{Pistes d'amélioration}
% ═══════════════════════════════════════════════════════════════════════════════

\begin{enumerate}[nosep]
  \item \textbf{Données supplémentaires} : le dataset (1\,682 sessions) reste limité. Plus de données aideraient le Transformer.
  \item \textbf{Features météo} : intégrer des données radar (réflectivité, vent) comme features additionnelles.
  \item \textbf{Recalibration} : appliquer une recalibration isotonique pour corriger la sur-couverture.
  \item \textbf{Hybride XGBoost + Hawkes} : combiner la précision tabulaire du XGBoost avec la composante séquentielle du Hawkes.
  \item \textbf{Encodage temporel appris} : remplacer le positional encoding fixe par un MLP appris.
\end{enumerate}


% ═══════════════════════════════════════════════════════════════════════════════
\section{Phase 2 ter : Spatial Mixture of Experts}
% ═══════════════════════════════════════════════════════════════════════════════

\subsection{Motivation spatiale}

Bien que le modèle GRU Gaussian offre d'excellentes performances globales, il présente un comportement typique en fin de session : il sous-estime souvent la durée au début puis la surestime systématiquement vers la fin de l'alerte.

Pour corriger ce phénomène, nous avons formulé l'hypothèse suivante : \textbf{la dimension spatiale, et particulièrement la cinématique de l'orage, est la clé pour prédire la fin imminente.} Si le système orageux, ou son centroïde, s'éloigne rapidement de l'aéroport, l'alerte devrait être levée plus tôt, même si l'intensité brute des éclairs reste élevée en bordure de radar.

\subsection{Nouvelles features spatiales}

Aux 12 features existantes ont été ajoutés 8 nouveaux descripteurs spatio-temporels :
\begin{itemize}[nosep]
  \item \textbf{Vitesse radiale (12)} : Taux de changement de la distance au capteur (dérivée temporelle), indiquant une approche ou un éloignement (km/min).
  \item \textbf{Dispersion angulaire (13)} : L'éclatement spatial de l'orage, calculé via l'écart-type de l'azimuth sur les 10 derniers événements.
  \item \textbf{Ratio 5 km (14)} : Proportion des éclairs récents frappant dans le rayon critique de 5 km.
  \item \textbf{Indicateur de repli (15)} : Signal binaire déclenché si la vitesse radiale indique un éloignement \emph{et} que le ratio 5 km baisse.
  \item \textbf{Durée de recul (16)} : Temps continu (en minutes) pendant lequel le système s'éloigne sans interruption.
  \item \textbf{Vitesse du centroïde (17-18)} : Composantes $v_x, v_y$ du mouvement du barycentre.
  \item \textbf{Distance minimale (19)} : Coup de foudre le plus proche dans la fenêtre des dernières 5 minutes.
\end{itemize}

\subsection{Architecture : Spatial Mixture of Experts (MoE)}

Pour traiter l'hétérogénéité des dynamiques spatiales, une architecture \textbf{Mixture of Experts} a été développée autour d'un encodeur Transformer.

\begin{figure}[H]
  \centering
  \begin{tcolorbox}[colback=white,colframe=meteoCyan,boxrule=1pt,arc=4pt]
    \centering
    \textsf{\small $\underbrace{\text{Features (20D)}}_{\downarrow} \xrightarrow{\text{Transformer (Shared)}} \text{Hidden State } h_t$}\\
    \vspace{0.5em}
    \textsf{\small $\underbrace{\text{Features Spatiales}}_{\downarrow} \quad + \quad h_t \xrightarrow{\text{Gating Network}} \text{Poids des experts } w_A, w_B, w_C$}\\
    \vspace{0.5em}
    \textsf{\small $h_t \xrightarrow{\text{Expert A (Approche)}} (\mu_A, \sigma_A)$}\\
    \textsf{\small $h_t \xrightarrow{\text{Expert B (Recul)}} (\mu_B, \sigma_B)$}\\
    \textsf{\small $h_t \xrightarrow{\text{Expert C (Stationnaire)}} (\mu_C, \sigma_C)$}\\
    \vspace{0.5em}
    \textsf{\small $\text{Combinaison : } \mu = \sum w_k \mu_k \quad , \quad \sigma^2 = \sum w_k(\sigma_k^2 + \mu_k^2) - \mu^2$}
  \end{tcolorbox}
  \caption{Architecture du modèle Spatial MoE avec gating guidé par la cinématique.}
\end{figure}

Le réseau comprend $N=3$ experts indépendants (Gaussiens, espace log). Le \emph{Gating Network} distribue le poids de la prédiction en utilisant directement les features spatiales (vitesse radiale, ratio 5km, indicateur de repli) comme \textbf{indices de contexte}.

Afin d'éviter l'effondrement sur un seul expert (mode collapse), une pénalité de \emph{Load Balancing} a été ajoutée à la fonction de perte combinée pour forcer l'entrainement homogène du Mixture of Experts.

\subsection{Résultats de l'évaluation}

Les résultats agrégés sur le jeu de test confirment l'apport de l'architecture :

\begin{table}[H]
\centering
\small
\begin{tabular}{@{}lccccc@{}}
\toprule
\textbf{Variante} & \textbf{MAE} & \textbf{RMSE} & \textbf{Médiane} & \textbf{Biais} & \textbf{P90} \\
\midrule
\textbf{MoE Spatial} (Phase 2ter) & 23.68 & 32.53 & 19.35 & \textbf{+4.62} & 45.16 \\
GRU Gaussian (Phase 2bis) & 21.33 & 30.25 & 15.96 & \textbf{+1.96} & 42.13 \\
\bottomrule
\end{tabular}
\caption{Performance globale du MoE Spatial par rapport au modèle GRU v3.}
\end{table}

A l'échelle macroscopique, le GRU global conserve la meilleure MAE globale, due à l'induction séquentielle forte du RNN.
Cependant, l'analyse des cas marginaux révèle l'intérêt qualitatif du MoE : \textbf{le modèle MoE Spatial a des intuitions de terminaison beaucoup plus nettes} lorsque l'orage modifie brutalement son cap. Lors de phases de recul prononcées, le MoE délègue le signal prédictif à son Expert Spécialisé, résultant en une annulation plus rapide et plus franche de la surestimation résiduelle du temps restant.

Le Tableau \ref{tab:moe-demo} détaille les sessions où le MoE surpasse significativement le modèle de base grâce à ce comportement d'expert.


% ═══════════════════════════════════════════════════════════════════════════════
\section{Conclusion}
% ═══════════════════════════════════════════════════════════════════════════════

\begin{keybox}{Résultats clés de la Phase 2bis}
\begin{itemize}[nosep]
  \item \textbf{MAE = 21.33 min} (vs 22.11 en Phase 2, amélioration de 3.5\%)
  \item \textbf{Biais = +1.96 min} (quasi nul)
  \item \textbf{Incertitude calibrée} à 96.4\% dans l'intervalle 2$\sigma$
  \item \textbf{Homogénéité entre aéroports} (écart max de 2.57 min)
\end{itemize}
\end{keybox}

Le \textbf{GRU Gaussian avec target logarithmique} est le grand gagnant. La leçon principale est que la formulation de la perte (Gaussian NLL en log-space) importe plus que l'architecture (Transformer vs GRU).

Le Neural Hawkes Gaussian est désormais le meilleur modèle tous critères confondus. Il surpasse nettement le XGBoost AFT (MAE~=~57.28 en Phase~2) tout en fournissant une estimation d'incertitude. Pour la Phase~3, c'est le candidat naturel pour la présentation interactive.

\end{document}
